\chapter{The Middle Class}

\section{Sanskritization: Modernizing or Traditionalizing?}

\begin{KeyConcept}
\begin{itemize}
\item \textbf{Sanskritization is both a modernizing and a traditionalizing force.}
\item It is a \textbf{modernizing force for the upper caste} --- who exchanged their accumulated cultural, economic and symbolic capital for new capital (modern western education, trade and commerce); the primary beneficiaries were the \textbf{Brahmins and Banias}.
\item It is a \textbf{traditionalizing force for the lower caste} --- who emulated the cultural lifestyle of the upper castes.
\item Contemporary dimension: reservation (an affirmative-action/westernizing force) uplifted the lower castes, which threatened the upper-caste position --- so the upper castes too began \textbf{demanding reservation} to retain their position (studied by \textbf{Harold Gould}).
\end{itemize}
\end{KeyConcept}

\section{Ramkrishna Mukherjee and the Rise of Class}

\begin{KeyConcept}
\begin{itemize}
\item \textbf{Ramkrishna Mukherjee} analysed 250 years of British rule through a \textbf{Marxist perspective} (like A. R. Desai), in two books: \textbf{'The Rise and Fall of the East India Company'} (as a historian) and \textbf{'The Dynamics of a Rural Society'} (as a sociologist --- a study of \textbf{Birbhum}, Bengal).
\item Before British rule the village was \textbf{egalitarian and undifferentiated} (a largely homogeneous body of peasants, with no class differentiation). With the British, \textbf{land became private property} --- and a \textbf{three-tier class structure emerged: landlords, peasants and landless agricultural labourers}.
\item This is \textbf{historical materialism}: class formation follows the change in the means of production. It provided the \textbf{base for later Marxist village studies} (Thorner's Malik/Kisan/Mazdoor; Gough's five classes).
\item \textbf{One class $\rightarrow$ three classes (differentiation) $\rightarrow$ five classes (more differentiation) = the formation of the middle class.}
\end{itemize}
\end{KeyConcept}

\section{Theories of Class}

\begin{KeyConcept}
\begin{itemize}
\item \textbf{Marx}: a class = people with a \textbf{similar relation to the means of production}.
\item \textbf{Weber}: a class = people with a \textbf{similar market situation}, hence equal life chances.
\item \textbf{Erik Olin Wright}: \textbf{contradictory class locations} --- one can be a member of two or three classes at once (a worker and a shareholder).
\item \textbf{Frank Parkin}: class is defined not by production/market but by the \textbf{strategy employed to exclude others and monopolise resources} (social closure).
\item \textbf{Anthony Giddens}: the \textbf{middle class is the educated class} (the entrepreneur = upper class; the worker = lower class).
\item \textbf{Shulamith Firestone}: \textbf{sex class} --- men and women are the historic classes.
\item \textbf{Pakulski \& Waters}: the \textbf{'death of class'} --- today \textbf{status (taste)} matters more than class.
\end{itemize}
\end{KeyConcept}

\section{Marx vs Weber on the Middle Class}

\begin{KeyConcept}
\begin{itemize}
\item \textbf{Marx} --- \textbf{polarization, not proliferation}: the middle class (the \textbf{petty bourgeoisie}) will \textbf{sink into the ranks of the proletariat} --- \textbf{proletarianization} (a term coined by \textbf{Raymond Aron}).
\item The laws of capitalism: the \textbf{law of general accumulation} (capitalists cut wages and install machines --- 'dead labour' replaces 'living labour' --- so workers lose purchasing power); then \textbf{overproduction + underconsumption $\rightarrow$ recession}, in which small businesses collapse and their owners \textbf{join the proletariat}. Hence there is no real scope for the middle class.
\item \textbf{Weber} --- \textbf{proliferation, not polarization}: he identified \textbf{four classes} (propertied upper class, propertyless white-collar workers, petty bourgeoisie, manual labourers); the \textbf{middle class (white-collar + petty bourgeois) will GROW} --- as the economy, consumption and \textbf{rationalization} (efficiency, calculability, predictability, control) expand.
\end{itemize}
\end{KeyConcept}

\section{Frank Parkin: Social Closure}

\begin{KeyConcept}
\begin{itemize}
\item \textbf{Frank Parkin's} strategies of social closure:
\item \textbf{Exclusionary social closure} --- the upper class excludes others and monopolises resources.
\item \textbf{Usurpationary social closure} --- the lower class 'usurps' resources (e.g. the Dalit movement --- temple entry, the Mahad Satyagraha --- seizing the ritual resource from the upper castes).
\item \textbf{Dual closure} --- a group that has usurped resources then acts like the elites and excludes its own: \textbf{Dalit men exclude Dalit women} (the Dalit movement is really a Dalit-men's movement). This is why \textbf{Dalit feminism} arose (Sharmila Rege, Jyotiba Phule, Savitribai Phule, Gail Omvedt) --- usurping from \textbf{upper-caste women and lower-caste men} (because the Indian women's movement was equated with upper-caste women's concerns).
\end{itemize}
\end{KeyConcept}

\section{The Indian Middle Class: B. B. Misra and Pavan Varma}

\begin{KeyConcept}
\begin{itemize}
\item \textbf{B. B. Misra} --- the \textbf{first scholar on the Indian middle class}: the origin of the contemporary middle class (doctors, engineers, bureaucrats, lawyers) can be traced to \textbf{colonialism and subsequent industrialization}. (Before the British, India had artisans, traders and merchants, but colonialism made this middle class \textbf{voluminous and visible}.)
\item \textbf{Pavan Varma} --- the middle class should \textbf{not be limited to consumption; it is equally about VALUES} (like \textbf{Goldthorpe's} 'affluent workers' --- who called themselves middle class for their new income, though they lacked the values: earning but reinvesting rather than merely consuming --- the \textbf{Protestant work ethic}; Benjamin Franklin's 'money begets money').
\item \textbf{Amartya Sen} (in '\textbf{The Argumentative Indian}') characterises the middle class as \textbf{argumentative} (never conceding a point) and self-appointed \textbf{custodian of Indian culture} --- but one must first be educated to argue.
\end{itemize}
\end{KeyConcept}

\section{Old vs New Middle Class}

\begin{KeyConcept}
\begin{itemize}
\item The \textbf{old middle class} is pre-1990s; the \textbf{new middle class} is post-\textbf{LPG reforms} (1990s onwards).
\item \textbf{What is 'new' about the new middle class}:
\item 1. It is the \textbf{post-90s consumption class} --- the main \textbf{driver of economic growth}.
\item 2. It is \textbf{more socially inclusive} --- not only upper castes (the old middle class --- doctors, lawyers, professors --- were largely Brahmins and Banias), though the low castes remain \textbf{underrepresented}.
\item 3. Its orientation is to the \textbf{private sector, IT and entrepreneurship} (start-up culture, 'unicorns') --- unlike the old middle class's orientation to government jobs; it looks to \textbf{Western education} (Oxford, Harvard, Silicon Valley) and treats \textbf{English as a passport to success}.
\item 4. It prefers \textbf{simple laws and hassle-free processes} (Digital India), dislikes subsidies, wants economic growth, and acts as a \textbf{watchful taxpayer} (RTI).
\item \textbf{Social composition} (2011): about \textbf{50\% of India} was in this 'new middle class', and \textbf{67\% of that 50\% were upper-caste Hindu men} --- the upper castes are overrepresented, the lower castes underrepresented.
\item The middle class is the \textbf{biggest vote bank} in India (influencers), and is concerned about environment, gender rights, water conservation, subsidies, and the tax-expenditure relation.
\end{itemize}
\end{KeyConcept}

\begin{QuickRevision}
\begin{itemize}
\item Sanskritization = modernizing (upper caste) + traditionalizing (lower caste); reservation $\rightarrow$ upper caste also demand reservation (Harold Gould).
\item Ramkrishna Mukherjee: egalitarian village $\rightarrow$ private property $\rightarrow$ three-tier class (landlord/peasant/landless labourer); base for Thorner \& Gough; 1 class $\rightarrow$ 3 $\rightarrow$ 5 = middle class.
\item Class theories: Marx (production), Weber (market), Wright (contradictory locations), Parkin (social closure), Giddens (educated = middle), Firestone (sex class), Pakulski-Waters (death of class).
\item Marx: polarization/proletarianization (middle class sinks); Weber: proliferation (middle class grows).
\item Parkin: exclusionary, usurpationary, dual closure (Dalit men exclude Dalit women $\rightarrow$ Dalit feminism).
\item Indian middle class: B. B. Misra (colonial origin), Pavan Varma (values, not just consumption; Goldthorpe).
\item Old (pre-90s) vs new (post-LPG): consumption driver, socially inclusive, private-sector/IT/entrepreneurial orientation, English = success, hassle-free/digital; 2011: 50\% new middle class, 67\% UCHM; biggest vote bank.
\end{itemize}
\end{QuickRevision}

